% Axial rod stretching — analytical derivation
% AIM Interactive Mechanics series. Compile with pdflatex.
\documentclass[10pt,a4paper]{article}
\usepackage[margin=2cm]{geometry}
\usepackage{amsmath}
\usepackage[T1]{fontenc}
\usepackage{lmodern}
\usepackage{microtype}
\usepackage{titlesec}
\titlespacing*{\section}{0pt}{9pt}{2pt}
\usepackage{xcolor}
\definecolor{iron}{HTML}{935636}
\definecolor{celdark}{HTML}{4B6C61}
\definecolor{celline}{HTML}{CDE0D8}
\usepackage[colorlinks=true,urlcolor=iron,linkcolor=black]{hyperref}
\setlength{\parindent}{0pt}
\setlength{\parskip}{4pt}
\AtBeginDocument{\setlength{\abovedisplayskip}{6pt}\setlength{\belowdisplayskip}{6pt}\setlength{\abovedisplayshortskip}{3pt}\setlength{\belowdisplayshortskip}{3pt}}
\pagestyle{empty}

\begin{document}

{\footnotesize\textcolor{celdark}{\textsc{aim interactive mechanics \,\textperiodcentered\, analytical derivation}}}

{\LARGE Axial Rod Stretching}

\textcolor{celdark}{\small Companion to the interactive simulator at
\url{https://aimlab.uk/sims/axial-stretching.html}}

\medskip
{\color{celline}\hrule height 1pt}

\section*{Setup and assumptions}
A straight rod of circular cross-section (diameter $d$) and length $L$ is fixed at
$x = 0$ and pulled by an axial force $F$ at $x = L$. The material is linear
elastic with Young's modulus $E$; strains are small. Let $u(x)$ be the axial
displacement of the cross-section at $x$.

\section*{Equilibrium}
Cut the rod at any $x$: axial equilibrium of the free piece requires the internal
normal force to equal the applied load, $N(x) = F$, at every section. Assuming the
stress is uniform over the cross-section,
\begin{equation}
\sigma \;=\; \frac{N}{A} \;=\; \frac{F}{A},
\qquad
A \;=\; \frac{\pi d^2}{4}.
\end{equation}

\section*{Kinematics and constitutive law}
The axial strain is the displacement gradient, and Hooke's law relates it to
stress:
\begin{equation}
\varepsilon \;=\; \frac{\mathrm{d}u}{\mathrm{d}x},
\qquad
\sigma \;=\; E\,\varepsilon .
\end{equation}

\section*{Governing equation}
Combining the three ingredients, $N = EA\,u'$, and equilibrium with no distributed
axial load gives $\mathrm{d}N/\mathrm{d}x = 0$, i.e.
\begin{equation}
EA \,\frac{\mathrm{d}^2 u}{\mathrm{d}x^2} \;=\; 0,
\qquad
u(0) = 0, \qquad EA\,u'(L) = F .
\end{equation}
This is the axial analogue of the fourth-order bending equation --- here the
problem is second order. Integrating twice: $u = C_1 x + C_2$; the boundary
conditions give $C_2 = 0$ and $C_1 = F/EA$.

\section*{Solution}
\begin{equation}
\boxed{\; u(x) \;=\; \frac{F x}{E A} \;}
\qquad\Rightarrow\qquad
\varepsilon \;=\; \frac{F}{EA} \;=\; \frac{\sigma}{E},
\qquad
\delta \;=\; u(L) \;=\; \frac{F L}{E A}.
\end{equation}
Strain and stress are uniform along the rod; the elongation is linear in load and
length, and inversely proportional to stiffness and area.

\section*{Worked example}
$E = 10$~MPa (a soft elastomer), $d = 10$~mm, $L = 500$~mm, $F = 50$~N:
\begin{equation}
A = \frac{\pi \times 10^2}{4} = 78.5~\text{mm}^2, \qquad
\sigma = \frac{50}{78.5\times10^{-6}} = 0.637~\text{MPa}, \qquad
\varepsilon = 6.37\%, \qquad
\delta = 31.8~\text{mm}.
\end{equation}
Set the simulator sliders to these values and compare. Then switch $E$ to steel
(200~GPa): the same load produces $\delta = 1.6~\mu$m --- the drawing does not
visibly move, and that is the point. Beyond roughly $10\%$ strain the
small-strain, linear-elastic assumptions degrade.

\vfill
{\color{celline}\hrule height 0.6pt}
\smallskip
{\footnotesize\textcolor{celdark}{Part of the multimodal mechanics series
(analytical \textperiodcentered\ simulated \textperiodcentered\ experimental)
\,\textperiodcentered\, AIM Group, Imperial College London.}}

\end{document}
